\documentclass[UTF8,a4paper,11pt,openany]{ctexbook}
\usepackage{../common/statlearnbook}
\SLSetBookMetadata{第 3 册}{统计学习方法讲义 v2.6.0}{优化模型与序列模型}
\begin{document}
\setcounter{page}{467}
\setcounter{chapter}{21}
\setcounter{figure}{0}
\pagestyle{headings}
\chaptermark{隐马尔可夫模型}
\begin{geometrybox}
同一状态格可支持三种递推：前向算法对进入同一结点的路径求和，后向算法对离开同一结点的路径求和，Viterbi 则取最大值并保存获胜前驱。
\end{geometrybox}
% v2.3.1 figure UID: FIG-P372-01
% Proposed post-v2.3.1 polish for FIG-P372-01: protect key-label size and stroke hierarchy.
\tikzset{slfig-FIG-P372-01/.style={font=\fontsize{9.2pt}{11.0pt}\selectfont,
  every path/.append style={line cap=round,line join=round}}}
\begin{figure}[H]
\centering
\begin{tikzpicture}[slfig-FIG-P372-01,
  state/.style={circle,draw=SLBlue,fill=SLBlueBG,line width=.78pt,minimum size=6.5mm,inner sep=0pt,font=\fontsize{9.2pt}{11pt}\selectfont},
  obs/.style={rounded corners=1mm,draw=SLRule,fill=SLSoftGray,line width=.60pt,minimum width=7mm,minimum height=5mm,font=\fontsize{9.2pt}{11pt}\selectfont},
  faint/.style={-{Stealth[length=1.5mm,width=.9mm]},draw=SLRule!78!SLTextGray,line width=.50pt},
  focus/.style={-{Stealth[length=2mm,width=1.2mm]},draw=SLBlue,line width=1.02pt},
]
\foreach \shift/\title/\kind in {0/{前向：汇总进入路径 $\sum$}/F,4.65/{后向：汇总离开路径 $\sum$}/B,9.30/{Viterbi：获胜路径 $\max$}/V}{
  \begin{scope}[xshift=\shift cm]
    \node[font=\bfseries\fontsize{9.5pt}{11.2pt}\selectfont,text=SLInk] at (1.15,2.50) {\title};
    \foreach \t/\lab in {0/{t-1},1/{t},2/{t+1}}{
      \node[font=\fontsize{8.7pt}{10.3pt}\selectfont,text=SLTextGray] at ({1.15*\t},2.05) {$\lab$};
      \node[state] (\kind\t a) at ({1.15*\t},1.18) {$q_1$};
      \node[state] (\kind\t b) at ({1.15*\t},.28) {$q_2$};
      \node[obs] (\kind\t x) at ({1.15*\t},-.72) {$x_{\the\numexpr\t+1\relax}$};
      \draw[draw=SLRule!72!SLTextGray,line width=.46pt,densely dashed] (\kind\t x)--(\kind\t b);
    }
    \foreach \t/\u in {0/1,1/2}{
      \draw[faint] (\kind\t a)--(\kind\u a);
      \draw[faint] (\kind\t a)--(\kind\u b);
      \draw[faint] (\kind\t b)--(\kind\u a);
      \draw[faint] (\kind\t b)--(\kind\u b);
    }
  \end{scope}
}
% Forward: multiple paths entering q_2 at time t.
\draw[focus] (F0a)--(F1b);
\draw[focus] (F0b)--(F1b);
% Backward: multiple paths leaving q_1 at time t.
\draw[-{Stealth[length=2mm,width=1.2mm]},draw=SLTeal,line width=1.02pt] (B1a)--(B2a);
\draw[-{Stealth[length=2mm,width=1.2mm]},draw=SLTeal,line width=1.02pt] (B1a)--(B2b);
% Viterbi: one winning path plus small backtrace arrows.
\draw[-{Stealth[length=2mm,width=1.2mm]},draw=SLGold,line width=1.08pt] (V0b)--(V1a)--(V2b);
\draw[-{Stealth[length=1.8mm,width=1.1mm]},draw=SLGold,line width=.78pt,densely dashed]
  ([yshift=2mm]V2b.west) to[bend left=18] ([yshift=2mm]V1a.east);
\node[font=\fontsize{8.8pt}{10.4pt}\selectfont,text=SLTextGray] at (6.98,-1.35)
  {同一状态格、时间刻度与观测行；浅灰边为非重点转移};
\end{tikzpicture}
\caption{HMM状态格的三种读法：前向与后向对共享路径求和，Viterbi只保存获胜前驱}\label{fig:V3-C05-lattice}
\end{figure}

图\ref{fig:V3-C05-lattice}中的浅灰边给出共同底格；彩色重点边只改变聚合读法，不改变状态、时间列或观测结构。
\textbf{递推与回溯。}前向、后向保留总概率，Viterbi 保留最大值及其前驱，因此后者需要回溯信息。
\end{document}
